\documentclass[11pt,a4paper]{article}

\usepackage[T1]{fontenc}
\usepackage[utf8]{inputenc}
\usepackage[french]{babel}
\usepackage{geometry}
\usepackage{graphicx}
\usepackage{hyperref}
\usepackage{enumitem}
\usepackage{booktabs}
\usepackage{xcolor}
\usepackage{titlesec}

\geometry{margin=2.5cm}
\hypersetup{
  colorlinks=true,
  linkcolor=blue!60!black,
  urlcolor=blue!60!black
}

\definecolor{keynest}{RGB}{24,24,27}
\definecolor{accent}{RGB}{249,115,22}

\titleformat{\section}{\large\bfseries\color{keynest}}{\thesection}{1em}{}
\titleformat{\subsection}{\normalsize\bfseries}{\thesubsection}{1em}{}

\title{
  \vspace{1cm}
  {\Huge\bfseries KeyNest}\\[0.4cm]
  {\Large Documentation des fonctionnalités}\\[0.2cm]
  {\large Dashboard Administrateur}
  \vspace{0.5cm}
}

\author{Projet \textit{mon-distributeur}}
\date{\today}

\begin{document}

\maketitle
\thispagestyle{empty}
\newpage

\tableofcontents
\newpage

% ---------------------------------------------------------------------------
\section{Introduction}

KeyNest est une plateforme de gestion de \textbf{distributeurs automatiques de clés} (casiers connectés). Le \textbf{dashboard administrateur} (\texttt{/admin}) centralise la supervision du réseau national, la gestion des réservations, des loueurs (hôtes type Airbnb), et le pilotage opérationnel en temps réel.

\subsection{Stack technique}
\begin{itemize}[noitemsep]
  \item \textbf{Frontend :} Next.js 16 (App Router), React 19, TypeScript, Tailwind CSS 4, Framer Motion
  \item \textbf{Backend :} Server Actions, API Routes
  \item \textbf{Base de données :} Prisma 7 + SQLite / Turso (LibSQL)
  \item \textbf{Cartographie :} Leaflet / React-Leaflet
  \item \textbf{Déploiement :} Vercel
\end{itemize}

\subsection{Réseau couvert (données de démonstration)}
14 villes, 25 distributeurs, 168 casiers répartis sur l'ensemble du territoire français (Paris, Lyon, Marseille, Nice, Bordeaux, Toulouse, Nantes, Lille, Strasbourg, Rennes, Montpellier, Grenoble, Nancy, Reims).

% ---------------------------------------------------------------------------
\section{Fonctionnalités transverses}

Ces éléments sont disponibles sur l'ensemble ou une grande partie du dashboard.

\subsection{Navigation}
\begin{itemize}
  \item \textbf{Sidebar fixe} (desktop) avec 10 sections principales
  \item \textbf{Barre de navigation mobile} (accès rapide : vue d'ensemble, carte, réservations, live, tower, journal)
  \item Indicateur d'\textbf{occupation globale} en bas de sidebar (casiers occupés / total, villes, distributeurs)
\end{itemize}

\subsection{Sélecteur réseau (\textit{Network Picker})}
Filtrage contextuel par :
\begin{itemize}[noitemsep]
  \item Ville (toutes les villes ou une ville précise)
  \item Distributeur (tous ou un site spécifique)
\end{itemize}
Le filtre s'applique aux KPI, réservations, historique, calendrier, journal et live grid.

\subsection{Bandeau KPI (KPI Strip)}
Affiché sur les onglets : Vue d'ensemble, Carte, Réservations, Historique, Calendrier.
\begin{center}
\begin{tabular}{ll}
  \toprule
  \textbf{Indicateur} & \textbf{Description} \\
  \midrule
  Casiers occupés & Nombre de casiers actuellement réservés \\
  Arrivées & Réservations débutant aujourd'hui \\
  Départs & Réservations se terminant aujourd'hui \\
  Pass en attente & Pass valides non encore scannés \\
  \bottomrule
\end{tabular}
\end{center}

\subsection{Rafraîchissement automatique}
\begin{itemize}[noitemsep]
  \item Polling des données toutes les \textbf{20 secondes}
  \item Notification toast lors d'un \textbf{scan réussi} (ouverture de casier détectée)
  \item Indicateur de santé API (\texttt{/api/admin/health}) dans l'en-tête
\end{itemize}

\subsection{Thème clair / sombre}
Basculement via le bouton en en-tête. Préférence persistée dans le \texttt{localStorage} du navigateur.

\subsection{En-tête dashboard}
\begin{itemize}[noitemsep]
  \item Titre dynamique selon l'onglet actif
  \item Statut \textit{Live} (connexion au backend)
  \item Sélecteur ville / distributeur
\end{itemize}

% ---------------------------------------------------------------------------
\section{Vue d'ensemble}

\textbf{Route :} onglet \og Vue d'ensemble \fg{} | \textbf{Objectif :} pilotage hiérarchique du réseau de casiers.

\subsection{Vue réseau national}
\begin{itemize}
  \item \textbf{Hero réseau} : taux d'occupation, barre de progression, 4 indicateurs (villes, distributeurs actifs/hors ligne, sites saturés)
  \item \textbf{Organisation par région} : accordéons par région administrative
  \item \textbf{Recherche} par nom de ville ou région
  \item \textbf{Tri} : occupation décroissante / croissante, ordre alphabétique
  \item Boutons \og Tout ouvrir \fg{} / \og Tout fermer \fg{}
  \item Régions à forte occupation ou avec sites offline \textbf{ouvertes par défaut}
  \item Code couleur : vert ($<$60\,\%), orange (60--85\,\%), rouge ($>$85\,\%)
  \item \textbf{Fil d'Ariane} : Réseau national $\rightarrow$ Ville $\rightarrow$ Distributeur
  \item Colonne latérale \textbf{Activité récente} (10 derniers événements)
\end{itemize}

\subsection{Vue par ville}
Cartes distributeurs avec : nom, adresse, taux d'occupation, barre de progression, badge saturation, statut hors ligne.

\subsection{Vue par distributeur}
\begin{itemize}
  \item En-tête site (adresse, occupation, statut offline)
  \item \textbf{Grille des casiers} : carte par casier (libre / occupé)
  \item \textbf{Ouverture à distance} sur casier occupé (bouton radio)
  \item Affichage code masqué, badge arrivée/départ, initiales du pass
  \item Activité récente filtrée par distributeur
\end{itemize}

% ---------------------------------------------------------------------------
\section{Carte du réseau}

\textbf{Objectif :} visualisation géographique et détection d'anomalies.

\subsection{Carte interactive (Leaflet)}
\begin{itemize}
  \item Carte de France avec marqueurs personnalisés (icônes clé / pin)
  \item Vue par \textbf{département} ou vue \textbf{France entière}
  \item Zoom sur département au clic
  \item Liste des distributeurs par département avec statistiques d'occupation
  \item Clic sur un distributeur $\rightarrow$ redirection vers la vue d'ensemble filtrée
\end{itemize}

\subsection{Radar d'anomalies}
Système de détection automatique en temps réel :

\begin{center}
\begin{tabular}{lll}
  \toprule
  \textbf{Type} & \textbf{Label} & \textbf{Déclencheur} \\
  \midrule
  \texttt{BLOCKED\_LOCKER} & Casier bloqué & Réservation expirée non récupérée \\
  \texttt{SCAN\_REFUSED} & Scan refusé & Tentative scan invalide / hors validité \\
  \texttt{DISTRIBUTOR\_OFFLINE} & Hors ligne & Distributeur \texttt{isActive = false} \\
  \texttt{OCCUPANCY\_PEAK} & Saturation & Occupation $\geq$ 85\,\% (warning) ou 100\,\% (critique) \\
  \bottomrule
\end{tabular}
\end{center}

\subsection{Overlay anomalies}
\begin{itemize}
  \item Pilule \og Anomalies \fg{} repliable sur la carte
  \item Accordion par type d'anomalie
  \item Filtre par type (bloqué, refusé, offline, saturation)
  \item \textbf{Faisceau visuel} (beam) vers l'anomalie sélectionnée sur la carte
  \item Clic sur anomalie $\rightarrow$ focus carte + zoom distributeur
\end{itemize}

% ---------------------------------------------------------------------------
\section{Casiers Live Grid}

\textbf{Objectif :} monitoring temps réel de tous les casiers du réseau.

\subsection{Fonctionnalités}
\begin{itemize}
  \item Grille live de tous les casiers (toutes villes ou filtré)
  \item Statuts : \textbf{Libre}, \textbf{Occupé}, \textbf{Bloqué}
  \item Filtres : Tous / Libres / Occupés / Bloqués
  \item \textbf{Recherche} par code, ville, distributeur, numéro de casier
  \item Regroupement par distributeur (accordéon)
  \item Bandeau d'occupation réseau animé
  \item \textbf{Pulse visuel} sur casiers avec activité récente (scan)
  \item \textbf{Panneau latéral détail} : infos casier, réservation, ouverture à distance
  \item Lien depuis le journal d'activité vers le casier concerné
\end{itemize}

% ---------------------------------------------------------------------------
\section{Réservations}

\textbf{Objectif :} gestion des passes actifs et futurs.

\subsection{Fonctionnalités}
\begin{itemize}
  \item Tableau des réservations actives et futures
  \item \textbf{Recherche} par code ou numéro de casier
  \item \textbf{Filtres} : toutes, aujourd'hui, cette semaine, utilisées, en attente
  \item \textbf{Tri} : par date d'arrivée ou de départ
  \item \textbf{Export CSV} des réservations visibles
  \item Actions par ligne :
    \begin{itemize}[noitemsep]
      \item Copier le lien pass (\texttt{/pass?code=...})
      \item Supprimer la réservation
    \end{itemize}
  \item Badges visuels : arrivée aujourd'hui, départ demain
  \item Affichage masqué du code (sécurité)
\end{itemize}

% ---------------------------------------------------------------------------
\section{Historique}

\textbf{Objectif :} consulter les séjours terminés ou expirés.

\begin{itemize}[noitemsep]
  \item Liste des réservations passées (dates de validité expirées)
  \item Filtrage par réseau (ville / distributeur)
  \item Informations : code, casier, dates, statut scanné ou non
\end{itemize}

% ---------------------------------------------------------------------------
\section{Calendrier}

\textbf{Objectif :} vision mensuelle des réservations.

\begin{itemize}[noitemsep]
  \item Vue calendrier mensuelle
  \item Navigation mois précédent / suivant
  \item Réservations affichées par jour (ville, distributeur, casier)
  \item Filtrage par réseau sélectionné
\end{itemize}

% ---------------------------------------------------------------------------
\section{Créer une réservation}

\textbf{Objectif :} création manuelle d'un pass client.

\subsection{Formulaire}
\begin{itemize}[noitemsep]
  \item Code de réservation (unique)
  \item Distributeur (groupé par ville)
  \item Date/heure d'entrée (\texttt{in}) — format français
  \item Date/heure de sortie (\texttt{out})
  \item Numéro de casier
\end{itemize}

\subsection{Après création}
\begin{itemize}[noitemsep]
  \item Génération automatique du \textbf{lien pass} (\texttt{/pass?...})
  \item Bouton \og Copier le lien \fg{}
  \item Enregistrement dans le journal (\texttt{MANUAL\_CREATE})
  \item Toast de confirmation
\end{itemize}

% ---------------------------------------------------------------------------
\section{Journal d'activité}

\textbf{Objectif :} traçabilité complète des opérations.

\subsection{Types d'événements enregistrés}

\begin{center}
\small
\begin{tabular}{ll}
  \toprule
  \textbf{Type} & \textbf{Description} \\
  \midrule
  \texttt{SCAN} & Scan QR réussi au distributeur \\
  \texttt{SCAN\_FAILED} & Scan refusé (code invalide ou hors validité) \\
  \texttt{REMOTE\_OPEN} & Ouverture à distance depuis le dashboard \\
  \texttt{CREATE} & Pass créé via API \\
  \texttt{MANUAL\_CREATE} & Réservation créée manuellement \\
  \texttt{DELETE} & Suppression de réservation \\
  \texttt{SUBSCRIPTION\_CREATED} & Nouvel abonnement loueur \\
  \texttt{SUBSCRIPTION\_UPDATED} & Modification statut abonnement \\
  \texttt{KEY\_REGISTERED} & Clé de logement enregistrée \\
  \texttt{KEY\_DEPOSITED} & Clé déposée en casier \\
  \texttt{KEY\_PASS\_CREATED} & Pass client généré pour un locataire \\
  \bottomrule
\end{tabular}
\end{center}

\subsection{Affichage}
\begin{itemize}[noitemsep]
  \item 100 derniers événements (polling)
  \item Filtrage par ville / distributeur (sélecteur réseau)
  \item Icônes et couleurs par type d'événement
  \item Mode compact disponible (utilisé dans vue d'ensemble et live grid)
\end{itemize}

% ---------------------------------------------------------------------------
\section{Loueurs \& Abonnements}

\textbf{Objectif :} gestion des hôtes (loueurs type Airbnb) et de leurs clés.

\subsection{Modèle économique — Formules}

\begin{center}
\begin{tabular}{lrrr}
  \toprule
  \textbf{Formule} & \textbf{Prix/mois} & \textbf{Clés max} & \textbf{Cible} \\
  \midrule
  Starter & 29\,\EUR{} & 3 & 1 à 3 logements \\
  Pro & 79\,\EUR{} & 10 & Loueurs actifs \\
  Business & 149\,\EUR{} & 25 & Parc immobilier étendu \\
  \bottomrule
\end{tabular}
\end{center}

\subsection{KPI loueurs}
\begin{itemize}[noitemsep]
  \item Abonnements actifs (+ essais)
  \item Revenu mensuel récurrent (MRR) estimé
  \item Clés utilisées / quota total
  \item Clés en casier (prêtes pour retrait client)
\end{itemize}

\subsection{Gestion des loueurs}
\begin{itemize}
  \item Création loueur : nom, e-mail, téléphone, société, formule
  \item Recherche et filtre par statut (actif, essai, impayé, résilié, sans abonnement)
  \item Fiche loueur extensible avec barre de quota clés
  \item Changement statut abonnement : Actif / Essai / Impayé / Résilié
\end{itemize}

\subsection{Gestion des clés (logements)}
\begin{itemize}
  \item Ajout clé : nom logement, adresse (sous quota abonnement)
  \item Statuts clé :
    \begin{itemize}[noitemsep]
      \item \textbf{À déposer} — enregistrée, pas encore en casier
      \item \textbf{En casier} — clé physiquement déposée
      \item \textbf{Récupérée} — retirée par un client
      \item \textbf{Inactive} — archivée (abonnement résilié)
    \end{itemize}
  \item Dépôt en casier : sélection distributeur + numéro de casier (vérification disponibilité)
  \item \textbf{Génération pass client} : code + dates $\rightarrow$ lien \texttt{/pass} pour envoi Airbnb
  \item Copie du lien pass en un clic
\end{itemize}

\subsection{Règles métier}
\begin{itemize}[noitemsep]
  \item Quota de clés selon formule (impossible de dépasser)
  \item Abonnement impayé/résilié : pas de nouvelles clés ni passes
  \item Vérification casier libre avant dépôt
  \item Lien réservation $\leftrightarrow$ clé loueur (\texttt{keyDepositId})
\end{itemize}

% ---------------------------------------------------------------------------
\section{Control Tower}

\textbf{Objectif :} centre de commandement opérationnel (sans carte ni grille).

\subsection{Health Score}
Score de santé réseau (0--100) avec label : Réseau sain / Sous surveillance / Dégradé / Critique.

\subsection{Morning Brief}
Résumé matinal automatique : anomalies, occupation, actions prioritaires.

\subsection{File d'actions (Kanban)}
Actions générées depuis les anomalies réseau :
\begin{itemize}[noitemsep]
  \item Priorités : Critique ($<$1h), Urgent ($<$4h), Attention ($<$24h), Info
  \item Statuts : À faire / En cours / Terminé
  \item Assignation opérateur (localStorage)
  \item Filtre par priorité et par région
  \item \textbf{Mode crise} : focus sur actions critiques uniquement
\end{itemize}

\subsection{Santé par région}
Cartes région avec score, nombre de distributeurs, nombre d'incidents.

\subsection{Performance Strip}
6 KPI avec sparklines sur 7 jours :
\begin{itemize}[noitemsep]
  \item Disponibilité réseau
  \item MTTR (temps moyen de résolution)
  \item Réservations actives
  \item Rotation moyenne
  \item Occupation réseau
  \item Scans / 7 jours
\end{itemize}

% ---------------------------------------------------------------------------
\section{Actions serveur (API \& Server Actions)}

\subsection{Admin (\texttt{actions.ts})}
\begin{itemize}[noitemsep]
  \item \texttt{getCities}, \texttt{getDistributors}, \texttt{getReservations}, \texttt{getActivityLogs}
  \item \texttt{createReservationManual} — création pass
  \item \texttt{deleteReservation}
  \item \texttt{remoteOpenBox} — ouverture à distance + log
\end{itemize}

\subsection{Loueurs (\texttt{landlord-actions.ts})}
\begin{itemize}[noitemsep]
  \item \texttt{getSubscriptionPlans}, \texttt{getLandlords}
  \item \texttt{createLandlordWithSubscription}
  \item \texttt{updateSubscriptionStatus}
  \item \texttt{addKeyDeposit}, \texttt{depositKeyInLocker}
  \item \texttt{createClientPassFromKey}
\end{itemize}

\subsection{API Routes}
\begin{itemize}[noitemsep]
  \item \texttt{POST /api/scan} — scan QR au distributeur (ouverture casier)
  \item \texttt{POST /api/pass} — création pass via API
  \item \texttt{POST /api/recuperer} — redirection lien collé (legacy)
  \item \texttt{GET /api/admin/health} — santé backend
\end{itemize}

% ---------------------------------------------------------------------------
\section{Pages publiques liées (écosystème voyageur)}

Hors dashboard admin, mais connectées au même système :

\begin{itemize}
  \item \textbf{\texttt{/pass}} — Affichage QR code + numéro de casier pour le voyageur
  \item \textbf{\texttt{/recuperer}} — Collage du lien Airbnb (espaces non encodés gérés) $\rightarrow$ redirection vers \texttt{/pass}
  \item \textbf{\texttt{/api/scan}} — Validation du scan et ouverture du casier
\end{itemize}

% ---------------------------------------------------------------------------
\section{Modèle de données (résumé)}

\begin{center}
\small
\begin{tabular}{ll}
  \toprule
  \textbf{Entité} & \textbf{Rôle} \\
  \midrule
  \texttt{City} & Ville du réseau \\
  \texttt{Distributor} & Distributeur / site de casiers (GPS, département) \\
  \texttt{Reservation} & Pass client (code, casier, dates, scanné) \\
  \texttt{ActivityLog} & Journal d'événements \\
  \texttt{SubscriptionPlan} & Formule d'abonnement loueur \\
  \texttt{Landlord} & Loueur / hôte \\
  \texttt{Subscription} & Abonnement actif d'un loueur \\
  \texttt{KeyDeposit} & Clé d'un logement en gestion \\
  \bottomrule
\end{tabular}
\end{center}

% ---------------------------------------------------------------------------
\section{Synthèse des 10 onglets}

\begin{center}
\small
\begin{tabular}{p{3.2cm}p{9.5cm}}
  \toprule
  \textbf{Onglet} & \textbf{Fonction principale} \\
  \midrule
  Vue d'ensemble & Navigation hiérarchique réseau $\rightarrow$ ville $\rightarrow$ casiers \\
  Carte & Géolocalisation + radar anomalies \\
  Réservations & CRUD passes actifs/futurs, export CSV \\
  Historique & Séjours expirés \\
  Calendrier & Planning mensuel \\
  Créer & Nouveau pass manuel \\
  Journal & Audit trail complet \\
  Casiers Live & Monitoring temps réel \\
  Loueurs & Abonnements + clés + passes clients \\
  Control Tower & Pilotage opérationnel + KPI \\
  \bottomrule
\end{tabular}
\end{center}

\vfill
\begin{center}
  \small\textit{Document généré pour le projet KeyNest — Dashboard Admin v1.0}
\end{center}

\end{document}
